\subsection{A qualitative repair of the ten-point geometric comparison}

The following argument supplies the finite recognition and fusion
input at the ten-point Klein configuration.  It does not assert the
particular numerical recursion in historical A.4.34--A.4.35 or the
compiler bound in A.4.57.  In particular, ordinary fixed loci are
not identified with profinite centralizers without an intervening
finite-cover argument.

\begin{result}{E.176}{Proved}
{Finite fixed-component recognition and finite fusion}
Let \(X\) be a connected aspherical smooth manifold of real
dimension \(d\), with a smooth action of a finite group \(C\).
Suppose that \(P=\pi_1(X,x)\) is finitely generated and is an
iterated extension of finitely generated free groups of rank at
least two.  Let \(H=C_2^2\leq C\), and choose \(x\in X^H\).
Let
\[
 1\longrightarrow P\longrightarrow E
   =\pi_1^{\rm orb}([X/C],x)\longrightarrow C\longrightarrow1
\tag{E.176.1}
\]
be the actual path extension; the fixed point gives its named
subgroup \(H\).
For every finite quotient \(q:E\to Q\), after joining the label
map to \(C\), there is a finite-index subgroup \(V\leq P\),
characteristic in \(P\) and contained in \(\ker q\), such that
\[
 \operatorname{im}\bigl(C_{E/V}(H)\longrightarrow Q\bigr)
                       =q(C_E(H)).
\tag{E.176.2}
\]
There is such a refinement with the following additional
property.  Choose one point \(x_j\) in each component of \(X^H\),
and paths \(\gamma_j:x\to x_j\).  They give actual subgroups
\[
 s_j(h)=\bigl(h,\gamma_j*(h\gamma_j)^{-1}\bigr)\in E.
\tag{E.176.3}
\]
Every homomorphism \(\iota:H\to E/V\) whose composite to \(C\)
is the named inclusion has, in \(Q\), the form
\[
             \operatorname{pr}_Q\iota(h)
                       =q(p)\,q(s_j(h))\,q(p)^{-1}
              \qquad(h\in H)
\tag{E.176.4}
\]
for some \(p\in P\) and some component \(j\).
The same refinement works for every such \(\iota\).
Only finitely many components will be used in E.177.
\end{result}

\begin{proof}
An element of the extension E.176.1 is a label \(c\in C\)
together with a homotopy class of paths from \(x\) to \(cx\).
Multiplication concatenates the first path with the \(c\)-image
of the second.  This description gives the action of \(E/U\)
on the connected cover \(Y_U=\widetilde X/U\) whenever \(U\)
is normal in \(E\).  It also proves E.176.3 directly, including
its homomorphism identity.

We first give the cohomology input, rather than assume
profinite goodness.  For every finite-index subgroup \(U\leq P\)
there is a finite-index subgroup \(U'\leq U\) such that
\[
 H^{>0}(U,\mathbb F_2)\longrightarrow H^{>0}(U',\mathbb F_2)
                         \quad\hbox{is zero}.
\tag{E.176.5}
\]
Here and below the cohomology in this construction is ordinary
discrete group cohomology.  For a free group \(F\) of rank \(r\geq2\)
put \(N=F'F^2\) and \(W=F/(N'N^2)\).  The group \(W\) is finite:
its order is \(2^{r+1+2^r(r-1)}\).  Its centre has trivial image
in \(F/N\).  To see this, use the regular \(C_2^r\)-cover of
the rose.  The squared \(j\)-th generator is the nonzero
mod-two cycle supported on the two \(j\)-edges based at
\(\{0,e_j\}\).  For every nonzero translation \(v\in C_2^r\)
choose \(j\) with \(v\notin\{0,e_j\}\).  That translation moves
the cycle to a distinct chain.  Graph homology embeds in
edge chains, so the deck action on \(H_1(N,\mathbb F_2)\)
is faithful.  A central element of \(W\) must act trivially
on that module, proving the assertion about its projection.

If \(F\) is normal in a group \(G\), the conjugation action
on \(W\) has a finite-index kernel \(J\) satisfying
\(J\cap F\subset F'F^2\).  Indeed the kernel of the inner
action of \(F\) on \(W\) is the inverse image of its centre.
Now induct on the length of the free extension tower.
For \(1\to F\to G\to B\to1\), take a subgroup \(B_0\leq B\)
as in the induction hypothesis and intersect its inverse
image with \(J\).  Write the resulting extension as
\(1\to F_1\to G_1\to B_1\to1\).  In the natural
Lyndon--Hochschild--Serre comparison only rows \(0,1\) occur.
The row-one coefficient map
\(H^1(F,\mathbb F_2)\to H^1(F_1,\mathbb F_2)\) is zero,
whereas the positive-degree row-zero map factors through
the zero restriction to \(B_0\).  Thus the map on each
associated graded of positive total degree is zero.
There are at most two filtration pieces.  Repeating the
construction inside \(G_1\) gives a composite that is zero
on the actual positive cohomology groups.  Finite-index
subgroups inherit the free extension tower, by intersection
and projection, and their free factors have finite ranks.
This proves E.176.5.

The subgroup in E.176.5 may be made characteristic in \(P\)
without losing the vanishing.  If \([P:U']=m\), its coset
action shows that
\[
 \kappa_m(P):=\bigcap_{\varphi:P\to\mathfrak S_m}\ker\varphi
                              \subset U'.
\tag{E.176.6}
\]
This intersection has finite index because \(P\) is finitely
generated.  It is characteristic, including under continuous
automorphisms of \(\widehat P\) after completion.  Start with
\(U_0=P\cap\ker q\) and repeat this construction to obtain
normal subgroups of \(E\)
\[
 U_0\supset U_1\supset\cdots\supset U_{(d+2)^2},
 \qquad H^{>0}(U_i,\mathbb F_2)
             \longrightarrow H^{>0}(U_{i+1},\mathbb F_2)=0
             \quad\hbox{as a map}.
\tag{E.176.7}
\]
The displayed zero denotes the restriction map, not a
claim that the target cohomology group is zero.
Let \(V=U_{(d+2)^2}\).

Here is the fixed-point argument with its convergence
and component conditions.  A smooth finite-group action
has an equivariant triangulation; for example see
\cite[opening Theorem, part 1, p.199; conventions, p.201]{LitIllman1978}.
One may include the chosen fixed point as a vertex and
barycentrically subdivide.  The cells have dimension at
most \(d\), and their stabilizers fix them pointwise.
Lift this triangulation to every \(Y_i=\widetilde X/U_i\).
It need not have finitely many cells.  Cellular cochains
then use products; a product of free-orbit cochain modules
for a finite cyclic group is coinduced and has zero Tate
cohomology, just as a finite free-orbit module does.
Thus no compactness of \(X\) or finiteness of the cell
set is being assumed here.

For \(A=C_2\), use its complete rank-one resolution with
alternating maps \(1-a,1+a\), and a cochain complex with
degrees between \(-1\) and \(d\).
The ordinary-cohomology-first spectral sequence has only
\(d+2\) terms in each total degree.  A map zero on ordinary
cohomology is zero on this page and strictly moves the
finite filtration on Tate hypercohomology.  A composite
of \(d+2\) such maps is zero on the abutment.
Restriction to the \(A\)-fixed subcomplex is a Tate
isomorphism: its kernel consists of the free-orbit
relative cellular cochains.  On the fixed subcomplex
the two horizontal maps are zero over \(\mathbb F_2\).
Consequently the zero Tate map gives zero on every
ordinary cohomology group of that fixed subcomplex.
The bounded vertical range justifies these comparisons
despite the complete resolution being unbounded in both
horizontal directions.

First use the canonical \(H\)-actions fixing the chosen
basepoint lifts \(y_i\).  Asphericity identifies the
cohomology restrictions in E.176.7 with those of \(Y_i\).
Since the covers are connected, the relative complexes
\(C^*(Y_i,y_i;\mathbb F_2)\) have zero maps on all
cohomology.  Apply the cyclic argument to a subgroup
\(A\leq H\), and then to \(H/A\) on the \(A\)-fixed
complexes.  After the \((d+2)^2\) maps it gives zero on
\[
 H^0(Y_0^H,\{y_0\};\mathbb F_2)
                  \longrightarrow H^0(Y_V^H,\{y_V\};\mathbb F_2).
\tag{E.176.8}
\]
The characteristic functions of components show that
every fixed point upstairs projects into the component
of \(y_0\) downstairs.  Fixed loci of smooth finite-group
actions are locally path connected, so points in that
component can be joined by a fixed path.

Let \(\bar g\in C_{E/V}(H)\).  Then \(\bar g y_V\) is
fixed, and its image in \(Y_0\) can be joined to \(y_0\)
by such a path.  Project that path to \(X\), and denote
it by \(\beta\).  Its endpoints are \(x,cx\), where \(c\)
is the label of \(g\).  Since the labels are retained,
\(c\) commutes with every \(h\in H\); moreover \(h\beta=\beta\).
Thus \((c,\beta)\in E\) centralizes \(H\) by the literal
path multiplication.  Its residue in \(E/U_0\), and
hence in \(Q\), equals that of \(g\), because the lifted
path has the prescribed endpoint.  This proves the
forward inclusion of E.176.2.  The reverse inclusion
is given by the images of ordinary commuting elements.

For fusion use the actions of an arbitrary \(\iota\)
and its projections on the same covers.  These actions
are equivariant under the covering maps, but no longer
fix the selected basepoints.  Replace the relative
complexes by
\[
 \widetilde C^*(Y_i)
      =\operatorname{Cone}\bigl(\mathbb F_2
                                 \longrightarrow C^*(Y_i;\mathbb F_2)\bigr),
\tag{E.176.9}
\]
with constants in degree \(-1\).  Connectedness and
E.176.7 again give zero cohomology maps in all degrees.
The same two cyclic Tate comparisons apply.
In the restriction to the fixed subcomplex the constants
cancel from the mapping cone, so its relative kernel
still consists exactly of free-orbit cellular cochains.
They imply a zero map on the reduced-cone cohomology
of the \(H\)-fixed loci after \((d+2)^2\) levels.
If \(Y_0^{\iota(H)}\) were empty, all the upper fixed
loci would be empty and the map in degree \(-1\)
would be the identity of \(\mathbb F_2\).  This
contradiction proves that \(Y_0^{\iota(H)}\) is nonempty.

Choose one of its fixed points \(z\), projecting to a
point of the component containing \(x_j\).  Join that
point to \(x_j\) within \(X^H\).  Its lift, preceded by
\(\gamma_j\), gives a lift fixed by \(s_j(H)\).
An ordinary deck element \(p\in P\) moves that lift to
\(z\).  The actions of \(\iota(h)\) and
\(p s_j(h)p^{-1}\) on \(Y_0\) cover the same map \(h\)
and fix the same point.  Uniqueness of a lift with a
given value makes them equal.  Passing to \(Q\)
proves E.176.4.  This proves fusion uniformly over
all the complements at the stated refined input.
\end{proof}
\begin{result}{E.177}{Proved}
{The actual ten-point Klein comparison and its permitted corrector}
Let \(H=\langle a,b\rangle=C_2^2\) act on the sphere by
\(a(z)=-z,\ b(z)=1/z\).  Mark the three exceptional pairs
\[
 A_\pm=(0,\infty),\qquad B_\pm=(1,-1),\qquad
 C_\pm=(i,-i)
\tag{E.177.1}
\]
and the regular orbit
\(D_h=h(z_0)\), \(h\in H\), for \(z_0=1+\sqrt2\).
Put \(M=\Gamma_{0,[10]}\), \(P=\Gamma_{0,10}\), and let
\(G\leq M\) be the inverse image of
\(C_{\mathfrak S_{10}}(H)\).  Write \(N=C_G(H)\).
For every finite quotient \(q:G\to Q\) there are finite
refinements recognizing \(q(N)\) as in E.176.2 and fusing
every elementwise labelled copy of \(H\) as in E.176.4.
In the fusion conclusion the correcting element can
always be chosen in
\[
                          P\ \cup\ P\theta ,
\tag{E.177.2}
\]
where \(\theta\) is one fixed half twist interchanging
only the labels \(C_+,C_-\).
It is killed by both four-point forgetting maps retaining
\(\{A_+,A_-,D_1,D_b\}\) and
\(\{B_+,B_-,D_1,D_a\}\).

Consequently
\[
              C_{\widehat G}(H)=\overline N .
\tag{E.177.3}
\]
Every continuous automorphism \(\alpha\) of
\(\widehat M\) that induces the identity on label
permutations admits one correcting element in
\(\widehat P\cup\widehat P\theta\) after which it fixes
the two named involutions pointwise and preserves
\(\overline N\).
These are qualitative finite and profinite comparisons;
no historical explicit refinement bound is asserted.
\end{result}

\begin{proof}
Use \(X=\mathcal M_{0,10}(\mathbb C)\).  Fixing three
labels identifies it with the ordered configuration
space of seven points on a triply punctured sphere.
Successive forgetting maps are locally trivial:
disjoint small discs about the retained points and
isotopies supported there supply local trivializations.
Their fibres are spheres with successively
\(3,4,\ldots,9\) points removed.  They are aspherical,
with free fundamental groups of ranks \(2,\ldots,8\).
Their homotopy exact sequences prove inductively that
\(X\) is aspherical and that \(P\) has the required
free extension tower.  Its real dimension is fourteen.
The label action is holomorphic.  Its actual orbifold
path extension is the stated subgroup \(G\) of the
mapping class group; compare the definitions and
orbifold interpretation in
\cite[\S\S1--2, pp.656--660]{LitSchneps2006Automorphisms}.
Thus E.176 applies, for example with \(16^2=256\)
cohomology-killing refinements.

Here is the exact fixed-locus calculation.  At a point
of \(X^H\), each labelled permutation \(h\) is realized
by a unique Möbius map, because three distinct marked
points determine that map.  These maps form a faithful
copy of \(H\).  Normalize the two \(A\)-labels to
\(0,\infty\) and the label \(B_+\) to \(1\).
The map \(a\) must be \(-z\).  The map \(b\) exchanges
\(0,\infty\), fixes \(1\), and is therefore \(1/z\).
The other \(B\)-label is \(-1\), and the two \(C\)-labels
are \(i,-i\) in one of their two possible orders.
All four \(D\)-labels are determined by \(D_1=z\).
The sole remaining condition is
\[
                  z\notin\{0,\infty,1,-1,i,-i\}.
\tag{E.177.4}
\]
Conversely every such \(z\), in either \(C\)-ordering,
gives a fixed configuration.  Hence \(X^H\) is
precisely two copies of the connected six-punctured
sphere.  No unlabelled identification of the two
components is made.

The centralizer of the label permutations is
\[
 C_{\mathfrak S_{10}}(H)
        =C_2[A]\times C_2[B]\times C_2[C]\times H[D].
\tag{E.177.5}
\]
The three size-two orbits have different stabilizers
in \(H\), so they cannot be interchanged by a
permutation commuting with \(H\).  Within them their
flips are independent; on the regular orbit the
centralizer is its group of translations.  This
proves E.177.5, of order thirty-two.
The single \(C\)-flip exchanges the two components
in E.177.4.  The fixed path subgroups E.176.3 can
therefore be represented by \(H\) and
\(\theta H\theta^{-1}\): an actual half twist supplies
the path from the first base configuration to the
one with its two \(C\)-labels interchanged.
E.176.4 then supplies a pure deck correction.
Taking its inverse when necessary gives exactly
the correcting cosets E.177.2.
The half twist is supported in a disc containing
the two \(C\)-marks and no others.  After those
marks are forgotten its mapping class is trivial,
which proves the two stated forgetting properties.

Equation E.176.2 at every finite output implies
E.177.3.  In fact its proof gives a slightly more
specific statement: the same closure is generated
densely by the ordinary path elements \((c,\beta)\)
with \(c\) commuting with \(H\) and \(\beta\) a path
inside the base component of \(X^H\).
This follows from the particular representative
constructed after E.176.8, and does not require a
separate ordinary centralizer classification theorem.

Finally \(\alpha(H)\) is an elementwise labelled
copy of \(H\).  At each finite output E.176 gives
a correction in E.177.2.  The corresponding closed
sets of corrections in
\(\widehat P\cup\widehat P\theta\) have the finite
intersection property: a common refinement handles
any finite set of outputs.  This space is compact.
Its intersection therefore supplies one correction
fixing both involutions at every finite level.
The corrected automorphism preserves their full
centralizer, which is \(\overline N\) by E.177.3.
This last compactness argument uses one fixed compact
correction space and a cofinal family of genuine
finite quotients.
\end{proof}
\begin{result}{E.178}{Proved}
{The two actual coloured forgettings have one common source}
With the ten-point marking of E.177, let
\[
 R_0=\langle v,d\mid(vd)^2=1\rangle,\qquad
 R=\widehat R_0,\qquad S(v)=d,\quad S(d)=v.
\tag{E.178.1}
\]
The two labelled four-point forgettings in E.177 induce
continuous surjections
\[
        q_b,q_a:\overline N/H\longrightarrow R,
                         \qquad q_a=S q_b .
\tag{E.178.2}
\]
The generators and basepoint comparisons can be chosen as
specified below so that the last equality is literal.
For arbitrary other fixed geometric basepaths it has one
retained inner factor, which can be absorbed in the
corresponding target comparison.
Both maps are restrictions of actual forgetting homomorphisms
on their finite-index label stabilizers, followed by the
effective coloured four-point quotient.
\end{result}

\begin{proof}
Let \(Z\) be the base component of \(X^H\) in E.177,
with coordinate \(z\) and with the \(C\)-ordering \((i,-i)\).
Write an element of E.177.5 by its three pair-flip bits
\((e_A,e_B,e_C)\) and its regular-orbit translation.
Normalization of the \(A\)-labels and \(B_+\) gives
the following literal transformations on the fixed locus:
the \(A\)-flip sends \((z,\delta)\) to \((z^{-1},-\delta)\),
the \(B\)-flip sends it to \((-z,-\delta)\), and the
\(C\)-flip sends it to \((z,-\delta)\), where
\(C_+=\delta i\).  The two regular-orbit translations
act on \(z\) by \(-z,z^{-1}\) and leave \(\delta\) fixed.
Thus the label subgroup preserving \(Z\) is
\[
 C^+=\{e_A+e_B+e_C=0\}\times H[D].
\tag{E.178.3}
\]
Its kernel on \(Z\) is exactly the named \(H\).
Its quotient acts freely on \(Z\) by
\(z\mapsto-z,z\mapsto1/z\), because all their fixed
points were removed.  The invariant coordinate is
\[
                 u=z^2+z^{-2},\qquad
 Z/(C^+/H)=\mathbb P^1_u\setminus\{\infty,2,-2\}.
\tag{E.178.4}
\]
The ordinary orbifold path group
\(E_H=\pi_1^{\rm orb}([Z/C^+],z_0)\) therefore has
quotient by its named \(H\) equal to the free
fundamental group of that triply punctured sphere.
One can also see the kernel directly:
\(C^+=H\times H[D]\), where \(H\) acts trivially
and \(H[D]\) acts freely on \(Z\).
The map \(E_H\to G\) induced by inclusion has image
\(N^{\rm fix}\subset N\).  Injectivity of this map
is not required.  The fixed-path conclusion in
E.177 proves
\[
                         \overline{N^{\rm fix}}=\overline N .
\tag{E.178.5}
\]

Let \(M_b\) preserve separately the two blocks
\(\{A_+,A_-\}\) and \(\{D_1,D_b\}\), and let \(M_a\)
preserve separately \(\{B_+,B_-\}\) and \(\{D_1,D_a\}\).
They are finite-index subgroups of \(M\).
Put \(N_i^{\rm fix}=N^{\rm fix}\cap M_i\).
The regular-orbit translation labels show
\[
 N^{\rm fix}=H N_i^{\rm fix},\qquad
 H\cap N_b^{\rm fix}=\langle b\rangle,\quad
 H\cap N_a^{\rm fix}=\langle a\rangle.
\tag{E.178.6}
\]
Indeed stabilizing a given two-element regular block
has index two, and the named \(H\) acts transitively
on the two such blocks.

The effective moduli orbifold of four points in two
unordered labelled pairs has coordinate
\(u=\lambda+\lambda^{-1}\), where
\(\lambda\in\mathbb P^1\setminus\{0,1,\infty\}\).
It has punctures at \(\infty,2\) and an order-two
cone at \(-2\).  Its orbifold fundamental group is
E.178.1.  The ineffective double flip of the two
pairs is killed in this quotient.  Thus the
forgotten image of \(b\), respectively \(a\), in
E.178.6 is killed.  The actual forgetting
homomorphisms on \(M_i\) consequently descend from
\(N_i^{\rm fix}\) to \(N^{\rm fix}/H\).

For the first quartet \(\{0,\infty,z,1/z\}\) the
coordinate is exactly E.178.4.  For the second
quartet use
\[
 r=\frac{z-1}{z+1},\qquad
        r^2+r^{-2}=\frac{2(u+6)}{u-2}=:T(u).
\tag{E.178.7}
\]
The Möbius map used here sends its \(B\)-pair
to \(0,\infty\) and its selected \(D\)-pair to
\(r,1/r\).  The transformation \(T\) exchanges
the two punctures and fixes the cone.  The
chosen \(z_0=1+\sqrt2\) has \(u_0=6\), fixed
by \(T\).  In the coordinate
\(w=(u-6)/(u+2)\), \(T\) is \(w\mapsto-w\);
the basepoint is \(0\), the cone is \(\infty\)
and the punctures are \(1,-1\).  Choose positive
cusp loops along the two radial arcs exchanged
by this rotation, with \(d\) at \(u=2\) and
\(v\) at \(u=\infty\).  This gives precisely
the based automorphism \(S\) in E.178.1.
Hence E.178.2 holds on the image of \(E_H\).
Each quotient map from the triply punctured
sphere to the orbifold is onto, so these maps
are surjective.

For completeness, all completion maps in this
argument use the ambient induced topology.
A finite-index subgroup \(M_i\leq M\) has its
full profinite completion as its closure.
Its actual forgetting map therefore extends
continuously to that closure.
Moreover
\[
 \overline N=H\,\overline{N_i^{\rm fix}},\qquad
 H\cap\overline{N_i^{\rm fix}}=H\cap N_i^{\rm fix}.
\tag{E.178.8}
\]
The first equality follows from E.178.5--E.178.6
and the finiteness of \(H\); the second follows
from the retained label quotient.  Since \(H\)
is central in \(\overline N\) and the appropriate
intersection in E.178.8 is killed, the extended
maps descend to \(\overline N/H\).  Compactness
and the ordinary surjectivity prove profinite
surjectivity.  The equality between the two
maps extends from the dense fixed-path image.
No unidentified completion of \(N/H\) is used.
\end{proof}

\begin{result}{E.179}{Proved}
{The full-GT mixed relation with its exact Schreier coefficient}
Let \(F=(\lambda,f)\) belong to the full profinite
Grothendieck--Teichmüller group.  In the index-two
subgroup of \(\widehat F(x,y)\) with \(x\) odd and
\(y\) even, define \(c=c_2(f)\) by
\[
 [f]_{\rm ab}=(0,c,-c)
       \quad\hbox{on the Schreier basis }
                       (x^2,y,xyx^{-1}).
\tag{E.179.1}
\]
Then the following identities hold in the actual
profinite orbifold and mixed braid groups:
\[
 \begin{aligned}
 f(d,v^2)d^{-4c}f(v,d^2)&=v^{4c}
                         \quad\text{in }R,\\
 f(u,w^4)&=w^{8c}f(u^2,w^2)u^{4c}(uw)^{-6c}
                         \quad\text{in }\widehat B_3 .
 \end{aligned}
\tag{E.179.2}
\]
The second identity takes place in the completion
of the actual index-three subgroup
\(\langle u,w^2\rangle\leq B_3\), with its retained
central generator \((uw)^3\).
This reconstructs the qualitative same-parameter
input of historical A.4.39 for full GT.  It does
not assert its explicit finite pentagon bound.
E.179 alone retains the Schreier coefficient.
The separate proof E.180 identifies it with the
historical symmetrizer coefficient; no additional
relation \((\mathrm{III}\prime)\) or later CM
covariance is concluded here.
\end{result}

\begin{proof}
The standard full-GT action on the mixed genus-zero
mapping class groups fixes label permutations and
commutes, up to the actual basepath inner comparisons,
with labelled forgetting.  These are the standard
genus-zero actions, not an assumed action on a
higher-genus tower; see
\cite[\S1, equation (1.1), author pp.2--3]
{SymNakamuraSchneps2000}.
Apply E.177 to its action on \(\widehat M\).
The single correcting element belongs to
\(\widehat P\cup\widehat P\theta\), hence to
both retained label stabilizers.  It fixes
the named \(H\) pointwise and preserves
\(\overline N\).  The two actual forgettings
of E.178 therefore give an induced
automorphism \(\alpha_N\) of \(\overline N/H\)
and two target inner comparisons.

Here is the target action with its marking made
explicit.  On \(B_3=\langle u,w\mid uwu=wuw\rangle\)
the standard formula is
\[
 u\longmapsto u^\lambda,\qquad
 w\longmapsto f(u^2,w^2)^{-1}w^\lambda f(u^2,w^2).
\tag{E.179.3}
\]
In the coloured subgroup put \(d=u,v=w^2\).
Its central cap has presentation \(R_0\).
By antisymmetry, if \(h=f(v,d^2)\), its raw
action sends \(d\) to \(d^\lambda\) and
\(v\) to \(h v^\lambda h^{-1}\).  Conjugating
by \(h^{-1}\) gives
\[
        \phi(v)=v^\lambda,\qquad
        \phi(d)=h^{-1}d^\lambda h .
\tag{E.179.4}
\]
This also follows directly from the four-strand
sphere formula: the actual projection sends
\((\sigma_1,\sigma_2,\sigma_3)\) to \((u,w,u)\);
its third conjugator becomes
\(f(u^2,u^{-2})=1\).
Thus no equivariance of an unspecified finite
centralizer quotient has been assumed.
The effective coloured quartet has exactly the
two positive cusp markings in E.178; fixed
geometric basepath changes merely change its
inner representative.

Consequently, for some \(g_b,g_a\in R\), the
actual two squares are
\[
 q_b\alpha_N=\operatorname{Inn}(g_b)\phi q_b,\qquad
 q_a\alpha_N=\operatorname{Inn}(g_a)\phi q_a .
\tag{E.179.5}
\]
Here one source automorphism \(\alpha_N\) occurs
in both equalities.  Substituting \(q_a=S q_b\)
and cancelling the surjection \(q_b\) proves
that \(\phi\) and \(\psi=S\phi S\) are inner
equivalent.  Any retained inner factor from
an alternative choice in E.178.2 is absorbed
in \(g_a\); it does not create two source
automorphisms or two independently chosen
Schreier coefficients.

We include the exact centralizer deduction.
In \(R_0=C_2*\mathbb Z\), write \(a=vd\).
For a finite quotient retaining \(a\), with
kernel \(K\), the associated graph has edges
the quotient group and vertices its cosets
modulo \(\langle a\rangle\).
If \(n\) is the order of the image of \(d\),
the cycle \(d^n\) has mod-two edge support
exactly the cyclic subgroup generated by
that image.  Its deck stabilizer is that
cyclic subgroup.  In the refinement
\(R_0/(K'K^2)\), every element centralizing
\(d\) must therefore project to a power of
\(d\) in the old quotient.  Conversely those
powers centralize \(d\).  As the old quotient
was arbitrary, this proves
\[
                    C_R(d)=\langle d\rangle^{\wedge}.
\tag{E.179.6}
\]
The same argument after \(S\) gives the result
for \(v\).  Unit powers have the same
centralizers.

Put \(k=f(d,v^2)\).  The equality of the two
inner-related actions first on \(d\), then
on \(v\), and E.179.6 give \(s,t\in\widehat{\mathbb Z}\)
with
\[
                              k d^t h=v^s .
\tag{E.179.7}
\]
Both \(h\) and \(k\) have zero abelianization
in \(R\).  Its abelianization is
\(\widehat{\mathbb Z}[d]\oplus C_2[a]\).
Thus \(t=-s\) and \(s\) is even.
The equality lies in the actual index-two
free subgroup
\[
                       K_{\rm std}=F(v,d^2).
\tag{E.179.8}
\]
Write its abelian basis as \(X=v,Y=d^2\).
The substitution \((x,y)\mapsto(d,v^2)\)
sends the Schreier basis in E.179.1 to
\[
                       Y,\quad X^2,\quad(X^{-1}Y^{-1})^2.
\tag{E.179.9}
\]
The last identity is the square of
\(dvd^{-1}=v^{-1}d^{-2}\).
It follows that \([k]=(4c,2c)\), whereas
\([h]=0\).  Abelianizing E.179.7 in
E.179.8 gives
\[
                      (4c,2c+t/2)=(s,0).
\tag{E.179.10}
\]
Division by two here is on the subgroup
\(2\widehat{\mathbb Z}\), where it is
unambiguous.  Hence \(t=-4c,s=4c\),
which proves the first identity in E.179.2.

Finally retain the centre.  The discrete
coloured braid subgroup has presentation
\[
 J=\langle d,V\mid(dV)^2=(Vd)^2\rangle,\qquad
              z=(dV)^2=(uw)^3,\quad d=u,\ V=w^2 .
\tag{E.179.11}
\]
Its central quotient is \(R_0\).
The inverse image of the free index-two
subgroup E.179.8 is a central extension of
a free group by \(\mathbb Z[z]\); it splits
because lifts of a free basis define a
section.  It is therefore \(F_2\times\mathbb Z\).
Taking full completions of this finite-index
subgroup proves that the kernel of
\(\widehat J\to R\) is precisely
\(\widehat{\mathbb Z}[z]\).
The first identity, using antisymmetry,
therefore gives
\[
               f(d,V^2)=V^{4c}f(d^2,V)d^{4c}z^e
\tag{E.179.12}
\]
for some \(e\in\widehat{\mathbb Z}\).
The actual character \(J\to\mathbb Z\)
sending \(d\mapsto1,V\mapsto0\) sends
\(z\) to \(2\).  Applying it to E.179.12
gives \(0=4c+2e\), hence \(e=-2c\).
This proves the second identity in E.179.2
with its central factor, not just its cap.

There is a useful precise literature boundary.
The Schreier coefficient is
\(-A_F(-1,1)\), by applying the two-state
Schreier abelianization to the commutator
generator; that generator has vector
\((0,-1,1)\), and conjugation by \(x\)
interchanges the last two coordinates.
Nakamura--Schneps identify this coefficient
with their symmetrizer-defined \(\rho_2\)
on the subgroup satisfying their additional
relation \(({\rm III}')\), not on all GT
merely by definition
\cite[Definition 5.1 and Proposition 5.3,
author pp.17--18]{SymNakamuraSchneps2000}.
The result here retains the exact
Schreier coefficient.  E.180 independently proves
its equality with the historical symmetrizer
coefficient by an explicit modular marking and
a continuous character.  Neither proof supplies
relation \((\mathrm{III}\prime)\) or the later CM
branch-covariance squares.
\end{proof}

